\inicializarSubseccion{Efecto de modificaciones de variables}

\tab Comparar el campo electrico generado por una carga puntual versus una carga distribuida en una linea recta cambia drásticamente nuestro ambiente. Es verdad que lo único que hicimos fue agregar mas cargas, pero el efecto que estas tienen es que nivelan en el centro el restos de las cargas cuando estas cargas estan distribuidas a lo largo de una linea recta. Pero si tenemos dos varas de distinta carga y de distinta longitud, el campo eléctrico generado por cada una de estas varas va a crear un campo no uniforme, haciendo posible la dielectroeforesis. De esta manera, podemos empezar a realizar simulaciones con campos eléctricos irregulares, lo cual es un gran paso para poder empezar a mover biomoleculas a través de campos eléctricos.